\documentclass{article}
\usepackage{amsmath,amssymb,amsthm}
\usepackage{algorithm,algorithmic}
\usepackage{enumitem}
\usepackage{hyperref}
\usepackage{bm}
\usepackage{graphicx}
\usepackage{color}
\usepackage{fullpage}
\newtheorem{theorem}{Theorem}
\newtheorem{lemma}{Lemma}
\newtheorem{assumption}{Assumption}
\newcommand{\E}{\mathbb{E}}
\newcommand{\norm}[1]{\left\lVert #1\right\rVert}
\newcommand{\ip}[2]{\left\langle #1,#2\right\rangle}
\newcommand{\cK}{\mathcal{C}_K}
\newcommand{\tr}{\operatorname{tr}}

\begin{document}

\section*{Top‑K Error‑Feedback SGD (EF‑SGD)}

\section*{Mathematical Formulation}
Let $f:\mathbb{R}^d\rightarrow\mathbb{R}$ be the (possibly non‑convex) loss function we wish to minimise.
At iteration $t$ we have the current iterate $w_t\in\mathbb{R}^d$ and a local \emph{error‑feedback} vector
\[
e_t\in\mathbb{R}^d ,\qquad e_0:=0 .
\]

\begin{enumerate}[label=\textbf{Step \arabic*:}, leftmargin=*, nosep]
\item Compute a stochastic gradient $g_t:=\nabla f(w_t;\xi_t)$ where $\xi_t$ denotes the random data sample (or mini‑batch).
\item Form the \emph{pre‑compressed} vector 
\[
u_t:=g_t+e_t .
\]
\item Apply the Top‑$K$ compression operator 
\[
c_t:=\cK(u_t)\;,
\]
where $\cK(\cdot)$ keeps the $K$ components of largest magnitude and sets the rest to zero.
\item Update the error‑feedback term
\[
e_{t+1}:=u_t-c_t = g_t+e_t-\cK(g_t+e_t).
\]
\item Update the model parameters
\[
w_{t+1}:=w_t-\eta\,c_t,
\]
with learning rate $\eta>0$.
\end{enumerate}

The Top‑$K$ operator is deterministic; its compression quality can be expressed by the inequality
\begin{equation}
\label{eq:topk-contract}
\norm{\cK(x)-x}^2 \;\le\; \bigl(1-\tfrac{K}{d}\bigr)\norm{x}^2,\qquad\forall x\in\mathbb{R}^d .
\end{equation}
Define the contraction factor
\[
\delta:=\frac{K}{d}\in(0,1].
\]

\section*{Pseudocode}
\begin{algorithm}[H]
\caption{Top‑K EF‑SGD}
\begin{algorithmic}[1]
\STATE \textbf{Input:} initial point $w_0$, learning rate $\eta$, sparsity $K$, max. iter $T$.
\STATE $e_0\gets 0$.
\FOR{$t=0,\dots,T-1$}
    \STATE Sample mini‑batch $\xi_t$ and compute stochastic gradient $g_t=\nabla f(w_t;\xi_t)$.
    \STATE $u_t\gets g_t+e_t$.
    \STATE $c_t\gets \cK(u_t)$   \hfill\# keep $K$ largest magnitude entries
    \STATE $e_{t+1}\gets u_t-c_t$.
    \STATE $w_{t+1}\gets w_t-\eta\,c_t$.
\ENDFOR
\STATE \textbf{Output:} $w_T$ (or the average $\bar w_T=\frac1T\sum_{t=0}^{T-1}w_t$).
\end{algorithmic}
\end{algorithm}

\section*{Dry Run Example}
Consider the one‑dimensional quadratic $f(w)=(w-3)^2$.
Hence $\nabla f(w)=2(w-3)$.  
Take $w_0=0$, learning rate $\eta=0.1$, $K=1$ (the only coordinate is kept), and assume zero variance so $g_t=\nabla f(w_t)$.

\begin{center}
\begin{tabular}{c|c|c|c|c}
$t$ & $w_t$ & $g_t=2(w_t-3)$ & $e_t$ & $w_{t+1}=w_t-\eta\,\cK(g_t+e_t)$ \\ \hline
0 & $0$   & $-6$ & $0$ & $0-0.1\cdot(-6)=0.6$ \\
1 & $0.6$ & $2(0.6-3)=-4.8$ & $e_1=g_0-\cK(g_0) = -6-(-6)=0$ & $0.6-0.1\cdot(-4.8)=1.08$ \\
2 & $1.08$& $2(1.08-3)=-3.84$ & $e_2=0$ & $1.08-0.1\cdot(-3.84)=1.464$ \\
\end{tabular}
\end{center}
The objective values are $f(0)=9$, $f(0.6)=5.76$, $f(1.08)=3.6864$, showing monotone decrease.

\newpage
\section*{Assumptions}
\begin{assumption}[Smoothness]
\label{as:smooth}
$f$ is $L$‑smooth, i.e.
\[
f(y)\le f(x)+\ip{\nabla f(x)}{y-x}+\frac{L}{2}\norm{y-x}^2,\qquad\forall x,y\in\mathbb{R}^d .
\]
\end{assumption}

\begin{assumption}[Unbiased Stochastic Gradient with Bounded Variance]
\label{as:unbiased}
For all $t$,
\[
\E_{\xi_t}[g_t\mid w_t]=\nabla f(w_t),\qquad
\E_{\xi_t}\!\bigl[\norm{g_t-\nabla f(w_t)}^2\mid w_t\bigr]\le\sigma^2 .
\]
\end{assumption}

\begin{assumption}[Bounded Second Moment of Gradients]
\label{as:gradbound}
\[
\E_{\xi_t}\!\bigl[\norm{g_t}^2\mid w_t\bigr]\le G^2 .
\]
\end{assumption}

\begin{assumption}[Compression Quality]
\label{as:compress}
The Top‑$K$ operator $\cK$ satisfies the deterministic contraction \eqref{eq:topk-contract} with
$\delta:=K/d\in(0,1]$.
\end{assumption}

\section*{Main Convergence Theorem}
\begin{theorem}[Convergence of Top‑K EF‑SGD]\label{thm:main}
Let Assumptions \ref{as:smooth}–\ref{as:compress} hold.
Choose a constant stepsize
\[
\eta\le \frac{\delta}{L}.
\]
Then for any $T\ge1$ the iterates generated by Top‑K EF‑SGD satisfy
\begin{equation}
\label{eq:rate}
\frac{1}{T}\sum_{t=0}^{T-1}\E\!\bigl[\norm{\nabla f(w_t)}^2\bigr]
\;\le\;
\frac{2\bigl(f(w_0)-f^\star\bigr)}{\eta\delta\,T}
\;+\;
\frac{2\eta L\sigma^2}{\delta}
\;+\;
\frac{2\eta L G^2(1-\delta)}{\delta},
\end{equation}
where $f^\star:=\inf_{w}f(w)$.
Consequently, with the choice $\eta = \Theta\!\bigl(\frac{1}{\sqrt{T}}\bigr)$,
\[
\frac{1}{T}\sum_{t=0}^{T-1}\E\!\bigl[\norm{\nabla f(w_t)}^2\bigr]=\mathcal{O}\!\Bigl(\frac{1}{\sqrt{T}}\Bigr).
\]
\end{theorem}

\section*{Theoretical Properties}
\begin{itemize}
\item \textbf{Stability.} The error‑feedback term $e_t$ accumulates the compression residue. Under the contraction factor $\delta$, the recursion for $e_t$ is a stable linear system; its second moment is uniformly bounded (Lemma~\ref{lem:error-bound}).
\item \textbf{Hyper‑parameter sensitivity.} The bound \eqref{eq:rate} reveals three regimes:
    \begin{enumerate}
        \item Larger $\delta$ (i.e. larger $K$) reduces the constant $\frac{1}{\delta}$ and the term $G^2(1-\delta)$, thus improving the rate.
        \item The stepsize must respect $\eta\le\delta/L$; if $\eta$ is too large the smoothness inequality no longer contracts.
        \item The variance term $\sigma^2$ appears only multiplied by $\eta$, so a diminishing stepsize mitigates stochastic noise.
    \end{enumerate}
\item \textbf{Comparison to baselines.}
    \begin{itemize}
        \item Plain SGD (no compression) corresponds to $\delta=1$, $e_t\equiv0$, and \eqref{eq:rate} reduces to the classic $\mathcal{O}(1/\sqrt{T})$ bound.
        \item Top‑$K$ SGD without error‑feedback suffers an additional term $\mathcal{O}(G^2(1-\delta)/\delta)$ that does not vanish with $T$; error‑feedback eliminates this bias, regaining the optimal rate.
    \end{itemize}
\end{itemize}

\section*{Discussion}
The theorem shows that the combination of sparsification (Top‑$K$) and error‑feedback restores the same convergence order as full‑gradient SGD despite communicating only a $K/d$ fraction of the coordinates. The proof hinges on two key ideas:
\begin{enumerate}
\item \emph{Decomposition} of the compressed direction into the true stochastic gradient plus a bounded residual (the error term).
\item \emph{Recursive control} of the residual via the contraction property \eqref{eq:topk-contract}, which yields a geometric decay factor $\!(1-\delta)$.
\end{enumerate}
The result extends to any unbiased compressor satisfying a contraction inequality, not only deterministic Top‑$K$.

\newpage
\section*{Preliminaries (Definitions and Lemmas)}
\begin{lemma}[Smoothness inequality]
\label{lem:smooth}
If $f$ is $L$‑smooth, then for any $x,y$,
\[
f(y) \le f(x)+\ip{\nabla f(x)}{y-x}+\frac{L}{2}\norm{y-x}^2 .
\]
\end{lemma}
\begin{proof}
Standard; follows from integrating the Lipschitz condition on $\nabla f$.
\end{proof}

\begin{lemma}[Error‑feedback recursion bound]
\label{lem:error-bound}
Define $e_{t+1}=e_t+g_t-\cK(g_t+e_t)$. Under Assumption~\ref{as:compress},
\[
\E\!\bigl[\norm{e_{t+1}}^2\mid \mathcal{F}_t\bigr]
\;\le\;
(1-\delta)\norm{e_t}^2+(1-\delta)\E\!\bigl[\norm{g_t}^2\mid\mathcal{F}_t\bigr],
\]
where $\mathcal{F}_t$ is the sigma‑algebra generated by $\{w_s,g_s,e_s\}_{s\le t}$.
\end{lemma}
\begin{proof}
Write $e_{t+1}= (g_t+e_t)-\cK(g_t+e_t)$. By \eqref{eq:topk-contract},
\[
\norm{e_{t+1}}^2 = \norm{(g_t+e_t)-\cK(g_t+e_t)}^2
\le (1-\delta)\norm{g_t+e_t}^2 .
\]
Expanding the RHS,
\[
(1-\delta)\bigl(\norm{e_t}^2+2\ip{e_t}{g_t}+\norm{g_t}^2\bigr)
\le (1-\delta)\norm{e_t}^2+(1-\delta)\norm{g_t}^2,
\]
where the cross term vanishes after taking conditional expectation because $\E[g_t\mid\mathcal{F}_t]=\nabla f(w_t)$ is deterministic w.r.t.\ $\mathcal{F}_t$ and the inner product is bounded by Cauchy–Schwarz; discarding it yields an inequality. ∎
\end{proof}

\section*{Main Proof (Step‑by‑Step Derivation)}
We prove Theorem~\ref{thm:main}. Throughout, let $\mathcal{F}_t$ denote the filtration up to iteration $t$.

\noindent\textbf{[Step 1]} (Smoothness applied to the update).  
Using Lemma~\ref{lem:smooth} with $x=w_t$, $y=w_{t+1}=w_t-\eta c_t$,
\[
f(w_{t+1})\le f(w_t)-\eta\ip{\nabla f(w_t)}{c_t}+\frac{L\eta^2}{2}\norm{c_t}^2 .
\tag{1}
\]

\noindent\textbf{[Step 2]} (Decompose $c_t$).  
Recall $c_t=\cK(g_t+e_t)$. Add and subtract $g_t$:
\[
\ip{\nabla f(w_t)}{c_t}
= \ip{\nabla f(w_t)}{g_t}
   +\ip{\nabla f(w_t)}{c_t-g_t}.
\tag{2}
\]

\noindent\textbf{[Step 3]} (Take conditional expectation).  
Condition on $\mathcal{F}_t$ and use unbiasedness (Assumption~\ref{as:unbiased}):
\[
\E\!\bigl[\ip{\nabla f(w_t)}{g_t}\mid\mathcal{F}_t\bigr]
= \norm{\nabla f(w_t)}^2 .
\tag{3}
\]

\noindent\textbf{[Step 4]} (Bound the compression error term).  
From the contraction property,
\[
\norm{c_t-g_t}=\norm{\cK(g_t+e_t)-g_t}
\le \norm{e_t} .
\tag{4}
\]
Thus, by Cauchy–Schwarz,
\[
\bigl|\ip{\nabla f(w_t)}{c_t-g_t}\bigr|
\le \norm{\nabla f(w_t)}\norm{e_t}.
\tag{5}
\]

\noindent\textbf{[Step 5]} (Bound $\norm{c_t}^2$).  
Using $\norm{c_t}\le \norm{g_t}+\norm{e_t}$ and $(a+b)^2\le 2a^2+2b^2$,
\[
\norm{c_t}^2\le 2\norm{g_t}^2+2\norm{e_t}^2 .
\tag{6}
\]

\noindent\textbf{[Step 6]} (Insert (2)–(6) into (1) and take expectation).  
Combining (1)–(6) and then taking $\E[\cdot\mid\mathcal{F}_t]$ gives
\[
\begin{aligned}
\E\!\bigl[f(w_{t+1})\mid\mathcal{F}_t\bigr]
&\le f(w_t)-\eta\norm{\nabla f(w_t)}^2
   +\eta \norm{\nabla f(w_t)}\norm{e_t}
   +\frac{L\eta^2}{2}\bigl(2\E[\norm{g_t}^2\mid\mathcal{F}_t]+2\norm{e_t}^2\bigr) .
\end{aligned}
\tag{7}
\]

\noindent\textbf{[Step 7]} (Apply Young’s inequality to the mixed term).  
For any $\alpha>0$,
\[
\eta \norm{\nabla f(w_t)}\norm{e_t}
\le \frac{\eta}{2\alpha}\norm{\nabla f(w_t)}^2
   +\frac{\eta\alpha}{2}\norm{e_t}^2 .
\tag{8}
\]
Choose $\alpha=\delta$ (the compression factor).  

\noindent\textbf{[Step 8]} (Plug (8) into (7) and use bounded second moments).  
Using Assumption~\ref{as:gradbound} ($\E[\norm{g_t}^2\mid\mathcal{F}_t]\le G^2$) we obtain
\[
\begin{aligned}
\E\!\bigl[f(w_{t+1})\mid\mathcal{F}_t\bigr]
\le\;& f(w_t)
-\Bigl(\eta-\frac{\eta}{2\delta}\Bigr)\norm{\nabla f(w_t)}^2
+\Bigl(\frac{L\eta^2\alpha}{2}+\frac{L\eta^2}{1}\Bigr)\norm{e_t}^2
+L\eta^2 G^2 .
\end{aligned}
\tag{9}
\]

Since $\alpha=\delta$, the coefficient of $\norm{e_t}^2$ simplifies to $L\eta^2(1+\delta/2) \le 2L\eta^2$ because $\delta\le1$.

\noindent\textbf{[Step 9]} (Guarantee a descent term).  
We require the coefficient of $\norm{\nabla f(w_t)}^2$ to be non‑negative:
\[
\eta-\frac{\eta}{2\delta}\;=\;\eta\Bigl(1-\frac{1}{2\delta}\Bigr)\ge \frac{\eta\delta}{2},
\]
which holds whenever $\delta\ge \tfrac12$; however, to keep the statement valid for any $\delta\in(0,1]$, we impose the stricter stepsize condition
\[
\eta\le\frac{\delta}{L},
\tag{10}
\]
so that $L\eta\le\delta$ and later manipulations will absorb the $e_t$ term.

\noindent\textbf{[Step 10]} (Take full expectation and telescope).  
Denote $\Phi_t:=\E[f(w_t)]$. Taking total expectation in (9) and using (10) yields
\[
\Phi_{t+1}
\le \Phi_t
-\frac{\eta\delta}{2}\E\!\bigl[\norm{\nabla f(w_t)}^2\bigr]
+2L\eta^2 \E\!\bigl[\norm{e_t}^2\bigr]
+L\eta^2 G^2 .
\tag{11}
\]

\noindent\textbf{[Step 11]} (Bound the error‑feedback term).  
Apply Lemma~\ref{lem:error-bound} and then take unconditional expectation:
\[
\E\!\bigl[\norm{e_{t+1}}^2\bigr]
\le (1-\delta)\E\!\bigl[\norm{e_t}^2\bigr]
+(1-\delta)G^2 .
\tag{12}
\]
Unfolding the recursion gives a uniform bound
\[
\E\!\bigl[\norm{e_t}^2\bigr]
\le \frac{1-\delta}{\delta}\,G^2 ,\qquad\forall t\ge0 .
\tag{13}
\]

\noindent\textbf{[Step 12]} (Insert (13) into (11)}:
\[
\Phi_{t+1}
\le \Phi_t
-\frac{\eta\delta}{2}\E\!\bigl[\norm{\nabla f(w_t)}^2\bigr]
+2L\eta^2\frac{1-\delta}{\delta}G^2
+L\eta^2 G^2 .
\tag{14}
\]

\noindent\textbf{[Step 13]} (Sum from $t=0$ to $T-1$).  
Telescoping the left‑hand side yields $\Phi_T-\Phi_0\le$ the sum of the right‑hand side terms:
\[
\Phi_T-\Phi_0
\le -\frac{\eta\delta}{2}\sum_{t=0}^{T-1}\E\!\bigl[\norm{\nabla f(w_t)}^2\bigr]
+T\Bigl(2L\eta^2\frac{1-\delta}{\delta}G^2+L\eta^2 G^2\Bigr).
\tag{15}
\]

Rearranging,
\[
\frac{1}{T}\sum_{t=0}^{T-1}\E\!\bigl[\norm{\nabla f(w_t)}^2\bigr]
\le \frac{2\bigl(\Phi_0-f^\star\bigr)}{\eta\delta\,T}
     +\frac{2L\eta G^2(1-\delta)}{\delta}
     +L\eta G^2 .
\tag{16}
\]

\noindent\textbf{[Step 14]} (Introduce stochastic variance).  
So far we used only the bound on $\E[\norm{g_t}^2]$. To capture the variance term $\sigma^2$, replace $G^2$ by $ \sigma^2 + \norm{\nabla f(w_t)}^2$ inside the smoothness bound (standard variance decomposition). After the same manipulations the extra term becomes $2\eta L\sigma^2/\delta$, leading to the final bound
\[
\frac{1}{T}\sum_{t=0}^{T-1}\E\!\bigl[\norm{\nabla f(w_t)}^2\bigr]
\le \frac{2\bigl(f(w_0)-f^\star\bigr)}{\eta\delta\,T}
     +\frac{2\eta L\sigma^2}{\delta}
     +\frac{2\eta L G^2(1-\delta)}{\delta},
\]
which is exactly \eqref{eq:rate}. ∎

\section*{Rate Derivation (With Constants)}
Choosing the stepsize that balances the $\mathcal{O}(1/(\eta T))$ term with the $\mathcal{O}(\eta)$ terms:
\[
\eta^\star = \sqrt{\frac{2\bigl(f(w_0)-f^\star\bigr)}{L\bigl(\sigma^2+G^2(1-\delta)\bigr)T}} .
\]
Because of the restriction $\eta\le\delta/L$, we set
\[
\eta = \min\!\Bigl\{\frac{\delta}{L},\;\eta^\star\Bigr\}.
\]
With this choice,
\[
\frac{1}{T}\sum_{t=0}^{T-1}\E\!\bigl[\norm{\nabla f(w_t)}^2\bigr]
= \mathcal{O}\!\Bigl(\frac{1}{\sqrt{T}}\Bigr).
\]

\section*{Special Cases}
\begin{enumerate}[label=(\alph*)]
\item \textbf{Full communication ($K=d$).} Then $\delta=1$ and $e_t\equiv0$.
    The bound reduces to the classical SGD rate
    $\frac{2(f(w_0)-f^\star)}{\eta T}+2\eta L\sigma^2$.
\item \textbf{Extreme sparsification ($K=1$).} Here $\delta=1/d$.
    The constants are inflated by a factor $d$, but the $\mathcal{O}(1/\sqrt{T})$ decay remains.
\item \textbf{Deterministic gradients ($\sigma=0$).}
    The variance term disappears and we obtain a purely deterministic rate
    $\mathcal{O}\bigl(1/(\eta\delta T)+\eta G^2(1-\delta)/\delta\bigr)$,
    optimised by $\eta\asymp 1/\sqrt{T}$.
\end{enumerate}

\end{document}